
SK 22.

From the primordial matter evolves the Great Principle;
from this evolves the I Principle;
from this evolves the set of sixteen;
from the five of this set of sixteen evolves the five elements.

SK 23.

Buddhi is ascertainment or will.
Virtue, knowledge, dispassion, and power are its manifestations
when the sattva attribute abounds.
And the reverse of these, when the tamas attribute abounds.

SK 24.

Ahamkara is self-assertion;
from that proceeds a two-fold evolution
the set of eleven and the five-fold primary elements

SK 25.

The set of eleven abounding in sattva proceeds from the Vaikriti form of I-Principle;
the set of five primary elements proceed from the Bhutadi form of I-Principle;
they are tamasa.
From the Taijasa form of I-Principle proceed both of them.

SK 26.

Organs of knowledge are called the Eye, the Ear, the Nose, the Tongue, the Skin.
Organs of action are called the Speech, the Hand, the Feet, the excretory organ and the organ of generation.

SK 27.

Of these (sense organs), the Mind possesses the nature of both
(the sensory and motor organs).
It is the deliberating principle, and is also called a sense organ
since it posseses properties common to the sense organs.
Its multifariousness and also its external diversities are owing
to special modifications of the Attributes.

SK 28.

The function of the five in respect to form and the rest,
is considered to be mere observation.
Speech, manipulation, locomotion, excretion and gratification
are the functions of the other five.

SK 29.

Of the three internal organs, their own characteristics are their functions;
this is peculiar to each.
The common modification of the instruments is the fives airs
such as prana and the rest.

SK 30.

Of all the four, the functions are said to be simultaneous and also successive
with regard to the seen objects;
with regard to the unseen objects, (and also seen objects)
the functions of the three are preceeded by that.

SK 31.

The organs enter into their respective being incited by mutual impulse.
The purpose of the Spirit is the sole motive (for the activities of the organs).
By none whatsoever is an organ made to act.

SK 32.

Organs are of thirteen kinds performing the functions of
seizing, sustaining and illuminating.
Its objects are of ten kinds, vis
the seized, the sustained, and the illumined.

SK 33.

The internal organ is three-fold.
The external is ten-fold; they are called the objects of the three (internal organs).
The external organs function at the present time and
the internal organs function at all the three times.
